% Section: The Evidence Request
%
% Promoted from a subsection of the diagnostics section. This is the operative
% handoff -- the form in which the frame leaves the paper and enters a letter,
% a procurement document, an examination procedure, or a complaint.

\section{The Evidence Request}
\label{sec:request}

The three diagnostics orient inquiry. They become shared work only when they
take the form of a specific request that a recipient can answer, refuse, or
answer partially with a stated reason. This section gives that form. Nothing in
it creates authority. It presumes a requester who already has standing to ask
--- an examiner under existing examination authority, a bank under its vendor
agreement, a regulator under its supervisory process, an affected party under
an applicable notice, complaint, appeal, or discovery procedure --- and
supplies only the structure of the asking.

The structure matters because the common failure is not refusal. It is a
responsive-looking answer that addresses a different proposition from the one
at issue. A request that names the proposition, the object it requires, and the
observations that would defeat it makes that substitution visible.

\subsection{The eight fields}
\label{sec:request:fields}

An evidence request derived from any of the three diagnostics should state
eight fields compactly. Stated as a request rather than a description:

\begin{enumerate}
  \item \textbf{Authority.} We are asking under [statute, regulation,
        enforceable condition, contractual provision, or adopted internal
        policy]. Supervisory guidance, where cited, informs the analysis and
        is not offered as the obligation: agency rules codifying the role of
        supervisory guidance provide that guidance does not carry the force
        and effect of law and is not by itself a basis for supervisory
        criticism. Identify any respect in which you contest the stated basis.
  \item \textbf{Proposition.} We are testing whether [the specific claim]. This
        is the proposition at issue; responses addressing a different claim
        should say so.
  \item \textbf{Required object.} Establishing that proposition requires
        evidence of [decision-time inputs and context / functional feature
        dependence / local behavioral approximation / internal computational
        state / a generated rationale / a human deliberative record /
        institutional authorization]. State which object you take the
        proposition to require if you disagree.
  \item \textbf{Offered artifact.} Identify the artifact you offer as
        establishing it, and the object that artifact exposes.
  \item \textbf{Binding records.} Provide the operative model and case binding,
        method and configuration, background-data construction, feature
        handling, software version, generation provenance, stability and
        relevant validation tests, known limitations, declared attestation
        scope, and reviewer access.
  \item \textbf{Defeating observations.} Identify the observations that would
        support the proposition and those that would defeat it. If no available
        observation could defeat it, say so.
  \item \textbf{Ownership.} Identify the party able to close or disclose any
        gap --- institution, vendor, or neither --- and the remedy available to
        the party whose interest the gap falls on.
  \item \textbf{Disposition.} State whether the result bears on a current
        finding, on an assessment of a claimed inference, or on a policy or
        design question outside the present review.
\end{enumerate}

Field~6 carries unusual weight. An artifact that no available observation could
have contradicted establishes that a process was run, not that its conclusion
holds. A validation regime with no failing condition, or a monitoring report
with no threshold that triggers action, answers fields~1 through~5 completely
and field~6 not at all.

\subsection{A worked request}
\label{sec:request:worked}

Consider FS~AI~RMF objective MS-2.9.1, which calls for structured model
explanations and method documentation, with validation for matters including
stability, performance, fairness, bias, and limitations. A bank offers a SHAP
report in support of the proposition that stated principal reasons are faithful
to a particular model output. The required object is method-specific functional
attribution under a declared value function and background construction; it is
not a transcript of model reasoning, and the report does not become one by
being called an explanation.

The request accordingly seeks the operative model and case binding, the value
function, the background-data construction, feature handling, software version,
generation provenance, stability and relevant validation tests, known
limitations, and reviewer access. Stable, reproducible attributions under the
declared construction, agreeing with suitable model-specific tests, support the
proposition. Instability, broken provenance, or material sensitivity to an
unjustified background defeats the broader claim while leaving the narrower one
intact.

The disposition then divides by actor. The bank can narrow its stated claim to
the bounded warrant the artifact supports, adopt a compensating control, or
disclose the residual gap. The vendor may be the only party able to supply the
missing provenance or export. An examiner may assess whether the bank's claimed
inference is supported, under authority the examiner already holds, without
prescribing any particular explanation architecture --- the objective does not
authorize a finding that one architecture was required. An affected party, where
a notice, complaint, appeal, or discovery process applies, can ask which method
produced the stated reasons and what binds that output to their own file, rather
than asking the institution to explain the denial.

That last division is the point of the form. The same eight fields, answered
once, tell four differently situated parties four different things they can each
act on separately.
